% Chapter 14: troubleshooting.
\chapter{Troubleshooting}\label{ch:trouble}
\index{troubleshooting}

Find the symptom in the left column. Try the actions in order. Switch the
power off before touching the paper path. If nothing helps, write down what
the screen shows and call [SUPPORT CONTACT].

\begingroup
\small
\renewcommand{\arraystretch}{1.4}
\begin{longtable}{@{}>{\raggedright\arraybackslash}p{0.25\textwidth}>{\raggedright\arraybackslash}p{0.30\textwidth}>{\raggedright\arraybackslash}p{0.38\textwidth}@{}}
\toprule
\textbf{Symptom} & \textbf{Likely cause} & \textbf{Action} \\
\midrule
\endfirsthead
\toprule
\textbf{Symptom} & \textbf{Likely cause} & \textbf{Action} \\
\midrule
\endhead
\midrule
\multicolumn{3}{r}{\emph{continued on next page}} \\
\endfoot
\bottomrule
\endlastfoot
%
\multicolumn{3}{@{}l}{\textbf{Buttons and starting}} \\
A push did nothing &
  The machine was not at the idle screen; it was just switched on; the push
  came too soon after another; or Test mode is on. &
  Check that the screen shows the idle screen (\shows{[1] 30yd [2] 90yd}).
  Wait one second after power-up. Leave a short moment between pushes.
  If it shows \shows{DEBUG: HOST CONTROL}, see below. If a push still does
  nothing, note which button and which screen, and report it. \\
A wind started by itself at power-up &
  Should not happen: a button already pushed in at power-up is ignored. &
  Switch off. Report it to [SUPPORT CONTACT]. \\
A button stays pushed in after a wind &
  Normal: the buttons are latching switches. &
  Nothing to do. The next push works whichever way it rests. \\
\btn{3} does nothing at idle &
  Normal: \btn{3} only pauses a running wind. &
  -- \\
\btn{4} does nothing during a wind &
  Normal: \btn{4} only works at idle or while paused. &
  Push \btn{3} to pause first. \\
Screen shows \shows{DEBUG: HOST CONTROL}, panel will not start &
  A PC has the machine in Test mode. &
  Finish or close the program on the PC, or push \btn{1} or \btn{2} to end
  Test mode. It also ends by itself 5 seconds after the PC stops talking. \\
%
\midrule
\multicolumn{3}{@{}l}{\textbf{During a wind}} \\
Paper jams, tears or bunches &
  Paper path problem. The machine does \emph{not} detect this during the main
  run. &
  Push \btn{1} or \btn{2} at once, or switch off. Clear the jam with the power
  off. Treat the roll as short. \\
Motor does not stop when \btn{1}/\btn{2} pushed &
  Controller or relay fault. &
  \textbf{Switch the power off.} Do not use the machine; call
  [SUPPORT CONTACT]. \\
Motor runs but the length stays at 0.0 &
  Wheel not on the paper, wheel slipping, or encoder cable loose. &
  Stop. Check the wheel and the encoder cable and connector. \\
\shows{ERR:n} on the top row, or \shows{!! ENC MISS:n !} &
  Encoder signals the controller could not read: loose connector, damaged
  cable, electrical noise. &
  Check the encoder cable and its connector. Keep it away from motor wires.
  Measure the next roll by hand. If the number keeps rising, call
  [SUPPORT CONTACT]. \\
Trim takes a long time with many bumps &
  Normal with the factory 2 yd coast when the real coast is much shorter (up
  to about 25 s); or the roll is heavy and slow. &
  Recalibrate (Chapter~\ref{ch:calibration}) to shorten the trim. \\
%
\midrule
\multicolumn{3}{@{}l}{\textbf{Results}} \\
\shows{!! UNDER TARGET !!} with \shows{SHORT} &
  Paper stopped moving during the trim (ran out, jam, wheel lifted), more than
  24 bumps were needed, or the trim took over 45 seconds. &
  Push a button to silence. Check paper, jam, wheel. Do not sell the roll as
  full length (Section~\ref{sec:under-target}). If it happens on good paper,
  recalibrate. \\
\shows{!! UNDER TARGET !!} with \shows{STOPPED} &
  You (or someone) pushed \btn{1} or \btn{2} before the end. &
  Normal. The roll is short by the amount shown. \\
Rolls are much longer than the label (many inches) &
  Coast setting too small, or wheel size wrong. &
  Recalibrate with the length check. \\
Rolls measure short by hand, but screen said \shows{DONE} &
  Wheel slipping, dirty or worn; wheel size wrong. &
  Clean the wheel, check it presses on the paper, recalibrate with the length
  check. Replace a worn wheel. \\
Result vanished from the screen &
  It clears itself 60 seconds after a good wind. A power cut or reset also
  clears it. &
  Normal for \shows{DONE}. After a power cut, measure the roll by hand. \\
%
\midrule
\multicolumn{3}{@{}l}{\textbf{Screen, LEDs and sound}} \\
Screen garbled or blank for a moment &
  Electrical noise from the relay. &
  Wait 2 seconds; it redraws itself. \\
Screen stays blank or garbled &
  Screen cable loose, or screen fault. &
  Switch off and on. Check the four-wire screen cable. \\
Screen blank and machine restarted by itself &
  The controller reset itself (watchdog) after a fault. The motor is stopped
  and the result is lost. &
  Measure the roll by hand. If it repeats, call [SUPPORT CONTACT]. \\
Fault LED lit (where fitted) &
  Internal controller fault. &
  Switch off and on. If it comes back, call [SUPPORT CONTACT]. \\
No beep at the end &
  You stopped the wind yourself (stops are silent), or buzzer fault. &
  Check the screen. If a full wind finished with \shows{DONE} and no beep,
  report the buzzer. \\
%
\midrule
\multicolumn{3}{@{}l}{\textbf{PC program}} \\
Program says no adapter, or cannot connect &
  Adapter not plugged in, missing driver, wrong wiring, no permission. &
  Run the \ui{Connection check} (Section~\ref{sec:doctor}) and follow its
  advice. Check TX to D0, RX to D1, GND to GND. \\
\ui{Firmware is compatible} fails &
  The firmware is too old or too new for this program. &
  Update the firmware (Chapter~\ref{ch:flashing}) or download the latest
  program. \\
\ui{Saved settings are valid} fails &
  Settings damaged, or firmware recently updated from before fw=6. &
  Calibrate and save once (Section~\ref{sec:settings-lost}). \\
Browser does not open &
  No default browser, or a locked-down computer. &
  Open \texttt{http://127.0.0.1:8765} yourself, or read
  \texttt{winder-open-me.txt}. \\
\ui{Stop motor} does not stop a wind &
  The wind was started from the panel. &
  Push \btn{1} or \btn{2} on the panel. \\
\end{longtable}
\endgroup
